\begin{figure}[htbp]
	\centering
	\resizebox{\textwidth}{!}{%
		\begin{tikzpicture}[
			trave/.style={line width=2.2pt, line cap=round, line join=round, black},
			terra/.style={line width=0.7pt, black},
			nodoV/.style={circle, fill=white, draw=black, inner sep=0pt,
				minimum size=5pt, line width=0.8pt},
			carico/.style={->, >=stealth, line width=1.0pt, black},
			reaz/.style={->, >=stealth, line width=1.0pt, gray!60},
			forza/.style={->, >=stealth, line width=1.0pt, black},
			blob/.style={draw=gray!45, fill=gray!10, line width=0.7pt,
				rounded corners=14pt},
			]
			
			\def\L{3.2}
			\def\sc{0.65}
			\def\rm{0.38}  % raggio archi momento — più grande per visibilità
			
			%% Macro arco ORARIO (~200° CCW, wraps RIGHT side)
			%% Visivamente: C aperta a sinistra, freccia punta in basso-sinistra in cima
			\newcommand{\arcORA}[1]{%
				\draw[->,>=stealth,line width=0.9pt]
				($(#1)+({\rm*cos(-100)},{\rm*sin(-100)})$) arc(-100:100:\rm);
			}
			%% Macro arco ANTIORARIO (~200° CCW, wraps LEFT side)
			%% Visivamente: C aperta a destra, freccia punta in basso-destra in fondo
			\newcommand{\arcANTI}[1]{%
				\draw[->,>=stealth,line width=0.9pt]
				($(#1)+({\rm*cos(80)},{\rm*sin(80)})$) arc(80:280:\rm);
			}
			%% Macro arco ANTIORARIO (~200° CCW, wraps LEFT side)
			%% Visivamente: C aperta a destra, freccia punta in basso-destra in fondo
			\newcommand{\arcANTIAP}[1]{%
				\draw[<-,>=stealth,line width=0.9pt]
				($(#1)+({\rm*cos(80)},{\rm*sin(80)})$) arc(80:280:\rm);
			}
			%% Macro arco ORARIO (~200° CCW, wraps RIGHT side)
			%% Visivamente: C aperta a sinistra, freccia punta in basso-sinistra in cima
			\newcommand{\arcORAAP}[1]{%
				\draw[<-,>=stealth,line width=0.9pt]
				($(#1)+({\rm*cos(-100)},{\rm*sin(-100)})$) arc(-100:100:\rm);
			}
			
			%% Macro arco ORARIO per asta VERTICALE
			%% (arcORA ruotato di 90°: avvolge in alto, aperto verso il basso,
			%%  attraversa l'asta verticale)
			\newcommand{\arcORAV}[1]{%
				\draw[->,>=stealth,line width=0.9pt]
				($(#1)+({\rm*cos(-10)},{\rm*sin(-10)})$) arc(-10:190:\rm);
			}
			
			%% Macro arco ANTIORARIO per asta VERTICALE
			%% (arcANTI ruotato di 90°: avvolge in basso, aperto verso l'alto,
			%%  attraversa l'asta verticale)
			\newcommand{\arcANTIV}[1]{%
				\draw[->,>=stealth,line width=0.9pt]
				($(#1)+({\rm*cos(170)},{\rm*sin(170)})$) arc(170:370:\rm);
			}
			
			%% Versioni con freccia capovolta (AP = freccia alla partenza)
			\newcommand{\arcORAVAP}[1]{%
				\draw[<-,>=stealth,line width=0.9pt]
				($(#1)+({\rm*cos(-10)},{\rm*sin(-10)})$) arc(-10:190:\rm);
			}
			\newcommand{\arcANTIVAP}[1]{%
				\draw[<-,>=stealth,line width=0.9pt]
				($(#1)+({\rm*cos(170)},{\rm*sin(170)})$) arc(170:370:\rm);
			}
			%% Layout globale:
			%% AB:  xshift=0,   B_AB  in x=3.2
			%% BC:  xshift=6.4, B_BC  in x=6.4   (gap tra B = 3.2)
			%% BD:  xshift=4.8  (midpoint 3.2+6.4/2=4.8), yshift=3.8
			%% DE:  xshift=6.4, yshift=3.8  (direttamente sopra BC)
			%%
			%% Blob BD:  da -1.6 a +1.3  → destra a 6.1
			%% Blob DE:  da -0.3 a +...  → sinistra a 6.1  ✓ nessuna sovrapposizione
			
			%% ════════════════════════════════════════
			%% TRATTO AB
			%% ════════════════════════════════════════
			\begin{scope}[xshift=0cm, yshift=0cm]
				\coordinate (A) at (0,0); \coordinate (B) at (\L,0);
				\coordinate (Mm) at (0.5*\L,0);
				\draw[blob](-0.80,-0.90) rectangle (\L+0.8, 0.90);
				\draw[trave](A)--(B);
				\draw[terra](A)--++(-0.22,-0.38)--++(0.44,0)--cycle;
				\draw[terra](-0.30,-0.38)--(0.30,-0.38);
				\foreach \x in {-0.24,-0.12,0,0.12,0.24}
				\draw[very thin,black](\x,-0.38)--++(-0.09,-0.11);
				\node[nodoV] at (A){};
				\draw[reaz](A)--++(0,\sc);
				\node[right,font=\small,gray!70] at ($(A)+(0.08,\sc*0.5)$){$p\ell$};
				\draw[carico](Mm)++(0,\sc)--(Mm);
				\node[left,font=\small] at ($(Mm)+(0,\sc)$){$p\ell$};
				\arcANTIAP{B}
				\node[above,font=\small] at ($(B)+(-0.50,0.3)$){$p\ell^2/2$};
				\node[left,font=\small] at (A){$A$};
				\node[right,font=\small] at (B){$B$};
			\end{scope}
			
			%% ════════════════════════════════════════
			%% TRATTO BC
			%% ════════════════════════════════════════
			\begin{scope}[xshift=6.5cm, yshift=0cm]
				\coordinate (B2) at (0,0); \coordinate (C) at (\L,0);
				\draw[blob](-0.80,-0.90) rectangle (\L+0.80, 0.90);
				\draw[trave](B2)--(C);
				\draw[terra](C)--++(-0.22,-0.32)--++(0.44,0)--cycle;
				\draw[terra]($(C)+(-0.30,-0.32)$)--($(C)+(0.30,-0.32)$);
				\draw[terra]($(C)+(-0.16,-0.44)$) circle(2.4pt);
				\draw[terra]($(C)+( 0.16,-0.44)$) circle(2.4pt);
				\draw[terra]($(C)+(-0.30,-0.56)$)--($(C)+(0.30,-0.56)$);
				\foreach \x in {-0.24,-0.12,0,0.12,0.24}
				\draw[very thin,black]($(C)+(\x,-0.56)$)--++(-0.09,-0.11);
				\node[nodoV] at (C){};
				\draw[reaz](C)--++(0,\sc);
				\node[left,font=\small,gray!70] at ($(C)+(-0.08,\sc*0.5)$){$p\ell$};
				\draw[forza](B2)++(0,\sc)--(B2);
				\node[right,font=\small] at ($(B2)+(0,\sc)$){$p\ell$};
				\arcORAAP{B2}
				\node[below,font=\small] at ($(B2)+(0.60,-0.2)$){$p\ell^2$};
				\node[left,font=\small] at (B2){$B$};
				\node[right,font=\small] at (C){$C$};
			\end{scope}
			
			%% ════════════════════════════════════════
			%% TRATTO BD  (verticale, centrato tra i due B)
			%% ════════════════════════════════════════
			\begin{scope}[xshift=4.8cm, yshift=1.8cm]
				\coordinate (B3) at (0,0); \coordinate (D) at (0,\L);
				\draw[blob](-1.10,-0.80) rectangle (1.0, \L+0.80);
				\draw[trave](B3)--(D);
				%% compressione B verso il basso
				\draw[forza](B3)--++(0,-\sc);
				\node[right,font=\small] at ($(B3)+(0.08,-\sc*0.5)$){$p\ell$};
				%% momento B ANTIORARIO: arco wraps LEFT, label a sinistra
				\arcORAV{B3}
				\node[left,font=\small] at ($(B3)+(-0.10,0.5)$){$p\ell^2/2$};
				%% compressione D verso l'alto
				\draw[forza](D)--++(0,\sc);
				\node[right,font=\small] at ($(D)+(0.08,\sc*0.5)$){$p\ell$};
				%% momento D ORARIO: arco wraps RIGHT, label a destra
				\arcANTIVAP{D}
				\node[right,font=\small] at ($(D)+(-1.20,-0.5)$){$p\ell^2/2$};
				\node[below left,font=\small] at (B3){$B$};
				\node[above left,font=\small] at (D){$D$};
			\end{scope}
			
			%% ════════════════════════════════════════
			%% TRATTO DE  (orizzontale, direttamente sopra BC)
			%% ════════════════════════════════════════
			\begin{scope}[xshift=6.5cm, yshift=5.0cm]
				\coordinate (D2) at (0,0); \coordinate (E) at (\L,0);
				\draw[blob](-0.55,-0.90) rectangle (\L+0.20, 1.20);
				\draw[trave](D2)--(E);
				\draw[black,line width=0.6pt]($(D2)+(0,0.65)$)--($(E)+(0,0.65)$);
				\draw[black,line width=0.6pt]($(D2)+(0,0.07)$)--($(E)+(0,0.07)$);
				\foreach \t in {0.00,0.17,0.34,0.50,0.67,0.84,1.00}
				\draw[carico,line width=0.7pt]
				($(D2)!\t!(E)+(0,0.65)$)--($(D2)!\t!(E)+(0,0.07)$);
				\node[above,font=\small] at ($(D2)!0.5!(E)+(0,0.65)$){$p$};
				%\draw[forza](D2)--++(0,\sc);
				%\node[left,font=\small] at ($(D2)+(-0.06,\sc*0.5)$){$p\ell$};
				\draw[forza](D2)--++(0,1.00);
				\node[left,font=\small] at ($(D2)+(0,1.40*0.5)$){$p\ell$};
				\arcORA{D2}
				\node[below,font=\small] at ($(D2)+(0.80,-0.10)$){$p\ell^2/2$};
				\node[below left,font=\small] at (D2){$D$};
				\node[below,font=\small] at (E){$E$};
			\end{scope}
			
		\end{tikzpicture}%
	}%
	\caption{Trave ramificata: isolamento dei quattro tratti come corpi liberi con le azioni alle connessioni interne.}
	\label{fig:corpi_liberi_ramificata}
\end{figure}
